\section{12个月共同承压规划测试}

以下冲击用于检验风险预算，不是历史回测、VaR、概率预测或最大损失承诺。每一资产层同时承受假设的不利变动；输入没有按历史频率或置信区间校准。交易成本、税费、跟踪误差、折溢价和再投资时点没有纳入。

\begin{exhibitbox}[基准稳健核心的共同承压算术]
\small
\begin{tabularx}{\exhibitboxwidth}{L{3cm}P P P}
\toprule
\textbf{资产层} & \textbf{权重} & \textbf{共同承压假设} & \textbf{加权贡献} \\
\midrule
组合内受保障流动性 & 20\% & 0\% & 0.00\% \\
交易现金 & 5\% & -1\% & -0.05\% \\
高质量人民币利率债 & 35\% & -5\% & -1.75\% \\
红利权益ETF & 30\% & -45\% & -13.50\% \\
非杠杆黄金ETF & 10\% & -10\% & -1.00\% \\
\midrule
\textbf{组合} & \textbf{100\%} &  & \textbf{-16.30\%} \\
\bottomrule
\end{tabularx}
\end{exhibitbox}

固定CNY200,000本金下，-16.30\%对应约CNY32,600。这一数字是共同承压情景的算术输出，用于讨论损失承受和流动性；它不会自行触发实际回撤复核，也不说明实际损失会停在该数值。

\section{实际回撤线与响应}

实际回撤必须以可追溯的账户价值记录测量：

\[
DD_t=\frac{\text{经外部现金流调整后的当期组合市值}}{\text{经外部现金流调整后的已记录高水位}}-1.
\]

每次记录至少应保留复核日期、账户或托管来源、当期市值、外部申购/取现、调整后高水位及数据质量。若这些字段缺失，不宣称已触发回撤线；情景表也不能替代该记录。

\begin{exhibitbox}[控制线]
\small
\begin{tabularx}{\exhibitboxwidth}{L{3cm}P P}
\toprule
\textbf{触发} & \textbf{阈值} & \textbf{研究响应} \\
\midrule
常规测量 & 每季度 & 记录权重、实际净分配、产品文件变化、折溢价、集中度和实际$DD_t$。 \\
再平衡复核 & 受保障流动性、利率债或总权益偏离目标5个百分点；交易现金超过CNY10,000或在分批期内低于目标；黄金偏离2个百分点 & 优先用新增资金和实际现金修复；不把下跌自动解释为买入信号。 \\
实际损失复核线 & 已记录$DD_t\leq-12\%$ & 暂停自动新增风险资产，重新核验资产层资格和用途匹配。 \\
实际风险处理线 & 已记录$DD_t\leq-18\%$ & 先检查流动性、杠杆、产品异常和权重，再由投资者作出非自动化决定。 \\
直接股票激活 & 合计不超过15\%，单一发行人不超过5\% & 只有发行人估值、现金流、税务、行业集中度和流动性审查全部通过。 \\
\bottomrule
\end{tabularx}
\end{exhibitbox}

\section{主要风险}

\subsection{利率、净值与流动性风险}

利率债ETF并非存款，利率上行时可发生净值下跌；货币ETF也不自动进入存款保险边界。任何产品在实施日前都需核验期限、费率、交易机制、折溢价和流动性。不能用“国债”或“货币”名称替代产品风险审查。

\subsection{权益因子和集中度风险}

红利权益ETF可以同时暴露于金融、周期和高分配行业。两只ETF并不天然分散，必须穿透看共同持仓、集团和行业权重。历史分配率下降、盈利变化或估值波动都可能使总回报和实际现金分配低于过去观察值。

\subsection{黄金与REITs风险}

黄金现货ETF不产生固定利息，且可能在共同承压中下跌。公募REITs有项目、估值和流动性风险；在三只独立项目未完成审查前，正确权重是0\%，而不是为了追求分配而提前买入。

\section{反方观点}

\textbf{质疑}：30\%红利权益ETF和35\%利率债仍可能同时下跌，3\%--6\%规划区间是否过于乐观？

\textbf{回应}：区间不是预测，也不用于证明策略“安全”。本案把共同承压、流动性缓冲、权重硬上限和实际回撤记录放在收益讨论之前；如果投资者无法接受约CNY32,600的规划压力损失，应选择更保守的向量或重新定义资金用途。
